%%%%%%%%%%%%%%%%%%%%%%%%%%%%%%%%%%%%%%%%%
% Journal Article
% LaTeX Template
% Version 1.3 (9/9/13)
%
% This template has been downloaded from:
% http://www.LaTeXTemplates.com
%
% Original author:
% Frits Wenneker (http://www.howtotex.com)
%
% License:
% CC BY-NC-SA 3.0 (http://creativecommons.org/licenses/by-nc-sa/3.0/)
%
%%%%%%%%%%%%%%%%%%%%%%%%%%%%%%%%%%%%%%%%%
%----------------------------------------------------------------------------------------
%       PACKAGES AND OTHER DOCUMENT CONFIGURATIONS
%----------------------------------------------------------------------------------------
\documentclass[paper=letter, fontsize=12pt]{article}
\usepackage[english]{babel} % English language/hyphenation
\usepackage{amsmath,amsfonts,amsthm} % Math packages
\usepackage[utf8]{inputenc}
\usepackage{float}
\usepackage{lipsum} % Package to generate dummy text throughout this template
\usepackage{blindtext}
\usepackage{graphicx} 
\usepackage{caption}
\usepackage{subcaption}
\usepackage[sc]{mathpazo} % Use the Palatino font
\usepackage[T1]{fontenc} % Use 8-bit encoding that has 256 glyphs
\linespread{1.30} % Line spacing - Palatino needs more space between lines
\usepackage{microtype} % Slightly tweak font spacing for aesthetics
\usepackage[hmarginratio=1:1,top=32mm,columnsep=20pt]{geometry} % Document margins
\usepackage{multicol} % Used for the two-column layout of the document
%\usepackage[hang, small,labelfont=bf,up,textfont=it,up]{caption} % Custom captions under/above floats in tables or figures
\usepackage{booktabs} % Horizontal rules in tables
\usepackage{float} % Required for tables and figures in the multi-column environment - they need to be placed in specific locations with the [H] (e.g. \begin{table}[H])
\usepackage{hyperref} % For hyperlinks in the PDF
\usepackage{lettrine} % The lettrine is the first enlarged letter at the beginning of the text
\usepackage{paralist} % Used for the compactitem environment which makes bullet points with less space between them
\usepackage{abstract} % Allows abstract customization
\renewcommand{\abstractnamefont}{\normalfont\bfseries} % Set the "Abstract" text to bold
\renewcommand{\abstracttextfont}{\normalfont\small\itshape} % Set the abstract itself to small italic text
\usepackage{titlesec} % Allows customization of titles

\renewcommand\thesection{\Roman{section}} % Roman numerals for the sections
\renewcommand\thesubsection{\Roman{subsection}} % Roman numerals for subsections

\titleformat{\section}[block]{\large\scshape\centering}{\thesection.}{1em}{} % Change the look of the section titles
\titleformat{\subsection}[block]{\large}{\thesubsection.}{1em}{} % Change the look of the section titles
\newcommand{\horrule}[1]{\rule{\linewidth}{#1}} % Create horizontal rule command with 1 argument of height
\usepackage{fancyhdr} % Headers and footers
\pagestyle{fancy} % All pages have headers and footers
\fancyhead{} % Blank out the default header
\fancyfoot{} % Blank out the default footer

\fancyhead[C]{} % Custom header text

\fancyfoot[RO,LE]{\thepage} % Custom footer text
%----------------------------------------------------------------------------------------
%       TITLE SECTION
%----------------------------------------------------------------------------------------
\title{\vspace{-15mm}\fontsize{24pt}{20pt}\selectfont\textbf{The EAGLE model with individual euro area countries: a framework for fiscal policy analysis}} % Article title
\author{
\large
{\textsc{Krzysztof~Ba\'nkowski\footnote{European Central Bank, Fiscal Policies Division - DG Economic Developements, Sonnemannstraße 20, 60314 Frankfurt am Main, E-Mail: krzysztof.bankowski@ecb.europa.eu}}}\\[2mm]
{\textsc{Emilė Petravičiūtė\footnote{European Central Bank, Fiscal Policies Division - DG Economic Developements, Sonnemannstraße 20, 60314 Frankfurt am Main, E-Mail: emile.petraviciute@ecb.europa.eu}}}\\[2mm]
%\thanks{A thank you or further information}\\ % Your name
%\normalsize \href{mailto:marco.torres.810@gmail.com}{marco.torres.810@gmail.com}\\[2mm] % Your email address
}
\date{}

%----------------------------------------------------------------------------------------
\begin{document}
\maketitle % Insert title
\thispagestyle{fancy} % All pages have headers and footers

% 1 Importance of heterogenity and interdependence
Interdependence and heterogeneity in the euro area are central to understanding the region’s macroeconomic dynamics, given its unique institutional framework, where a common monetary policy operates alongside decentralized fiscal and structural policies. Interdependence arises primarily through strong trade linkages and the transmission of monetary policy, implying that shocks in one member state—such as fiscal adjustments—can propagate across the union, amplifying volatility and influencing policy effectiveness. At the same time, persistent heterogeneity, reflected in disparities in productivity, debt dynamics, and structural rigidities, stems from differentiated policy priorities, institutional legacies, and varying degrees of economic resilience. These asymmetries pose challenges for union-wide stabilization efforts and risk exacerbating economic fragmentation. A comprehensive understanding of both interdependence and heterogeneity within the Economic and Monetary Union (EMU) is therefore essential for designing policies that mitigate fragmentation risks, enhance shock absorption, and reconcile national policy autonomy with collective stability. In this context, the development of analytical tools to improve euro area-wide institutions' capacity—such as the European Central Bank (ECB)—to assess and address these dynamics is of critical importance.


% 2 We establish a more granular version of EAGLE with 11 original member states
A key macroeconomic framework employed within the Eurosystem to analyze both interdependence and heterogeneity in the EMU has been the EAGLE model, developed in the 2010s (\cite[Gomes et al. (2012)]{Gomes}). This multi-country extension of the seminal New Area-Wide Model initially represented the euro area as a two-region economy—Germany and the rest of the euro area—alongside the United States and the rest of the world. While EAGLE has proven valuable for assessing economic policies within a multi-country setting, its structure limits its ability to evaluate large-scale initiatives such as the Next Generation EU (NGEU) program, which involves complex cross-country dimensions. Moreover, the model overlooks the role of smaller euro area economies, which, despite their size, contribute significantly to the system-wide dynamics of the union. To address these limitations, we build upon the EAGLE framework by further disaggregating the euro area into 11 early euro-adopting countries (Austria, Belgium, Finland, France, Germany, Greece, Ireland, Italy, Netherlands, Portugal, and Spain). This extension enhances the model's capacity to explicitly capture country-specific dynamics, providing a more granular representation of economic interdependencies within the EMU.

% 3 The model embdes a meaningful fiscal policy
Our model extension incorporates a meaningful role for fiscal policy, which was absent in the original EAGLE framework but introduced in later extensions (see \cite[Clancy et al. (2012)]{Clancy}). Specifically, government expenditure is disaggregated into government consumption and investment. Government consumption enters the utility function as part of a composite consumption basket alongside private consumption, thereby generating utility for households. In this setup, public consumption is not purely wasteful but exhibits complementarity with private consumption. Similarly, government investment brings benefits as it contributes to the accumulation of public capital, which enhances the productive capacity of the economy over time. Additionally, our model accounts for the direct import content of both types of government spending, recognizing that, a share of public purchases is allocated to imported goods, especially in small open economies.

% 4 Trade calibration is a major challange
A crucial aspect to be incorporated into our model extension is the significant role of trade in interconnecting the euro area and linking it to the global economy. The euro area is characterized by deep regional integration, facilitated by the single currency and the European Single Market, which has significantly reduced trade barriers. Intra-euro area trade constitutes a substantial share of member states' total trade flows, with Germany often serving as a central node, both as a major exporter and importer. At the same time, the euro area maintains a strong trade presence globally, ranking among the world's largest trading blocs. Its openness to trade, with total trade exceeding 30\% of GDP, underscores its importance as a global trading hub. When measured by the import-to-GDP ratio, euro area economies exhibit significantly higher openness compared to other major advanced economies, such as Japan and the United States (see Figure \ref{fig1}). Additionally, the euro area has historically maintained sizable trade surpluses, largely driven by its strong manufacturing sector. These trade dynamics are critical for understanding macroeconomic shock transmission within the euro area and beyond, and we aim to incorporate these features comprehensively in our extended model.

\begin{figure}[H]
    \centering
    \includegraphics[width=150mm]{/Users/kk/Documents/0000-00_work/2024-01_eagle/docs/2024-12_RCC-workshop/figures/ea_imports_chart.png}
    \caption{The magnitude of imports across euro area countries, compared to Japan and the US. \label{fig1}}
\end{figure}

% 5 This is how we calibrate trade


% 6(a) Having the model we study fiscal shocks (union wide)

% 6(b) National shocks

% 7(a) NGEU description

% 7(b) NGEU macro effects

% 7(c) NGEU spill-overs


% 5
Once known, the additional fiscal adjustment can be used in euro area macroeconomic simulations. The analysis of the second step makes use of the semi-structural model for the euro area ECB-BASE, which contains a rich specification of the general government. \footnote{For a description of the model, including its fiscal block, see \cite[Ba\'nkowski. (2023)]{Bankowski}} The actual fiscal policy is replaced in the model by the one consistent with full SGP compliance. The model is then re-simulated for the EMU period until 2019 and a counterfactual scenario is obtained. The simulations are conducted under the assumption of an exogenous monetary policy, notably with the policy rate remaining unchanged regardless of the fiscal policy course.

% 6
According to the simulations, full compliance with the SGP would have markedly changed the past macroeconomic situation of the euro area. The initial fiscal adjustment would have tamed the brisk output growth seen during the early years of EMU (see Figure \ref{fig2}). Significantly, the SGP would have ensured that sufficient buffers were in place, thus enabling an unfettered fiscal counter-response to the Great Financial Crisis. Furthermore, maintaining a broadly balanced structural position throughout the EMU period would have done away with the need to consolidate public finances. In reality, a major fiscal retrenchment took place amid the sovereign debt crisis and it aggravated the downturn, as illustrated by the simulations.

% 7
Adherence to the SGP would have left public finances in a significantly sounder shape than the state they were in. The maintenance of a structurally balanced fiscal position would have implied consistently higher headline balances, even with surpluses having been reached in good times. This would have kept the government debt ratio on a much more stable trajectory than what ultimately transpired.

\begin{figure}[H]
    \centering
    \includegraphics[width=150mm]{/Users/kk/Documents/0000-00_work/2021-08_SGP-reform/SGP_Reform/docu/note/figures/panel_EA_box2.pdf}
    \caption{Counterfactual scenarios assuming full SGP compliance with ECB-BASE (percentage of GDP except real GDP growth and GDP deflator growth, which are expressed in percentage year-on-year rates). \label{fig2}}
\end{figure}

Our findings indicate that non-compliance with the SGP framework significantly contributed to its failure in achieving its objectives. Notably, adherence to the SGP would have facilitated the building of fiscal buffers in the euro area and mitigated the tendency towards pro-cyclical fiscal policies. Ultimately, the effectiveness of any framework is contingent upon its implementation.

Our paper does not attempt to defend the pre-reformed SGP, as this would imply no need for reform of the governance framework. We acknowledge the numerous shortcomings of the SGP, as laid out in the recent ECB's Occasional Paper (see \cite[Haroutunian et al. (2024)]{Haroutunian}). Even with perfect implementation, the framework likely would not have ensured adequate government investment in the euro area. Additionally, it is improbable that it would have maintained a fiscal stance aligned with the monetary policy needs of the union. The reliance on unobservable indicators also likely contributed to non-compliance. Nevertheless, we emphasize that compliance issues contributed to the SGP's failure to achieve its objectives. This underscores a critical lesson for the future: the success of the newly reformed framework will depend on adherence to its provisions.

\begin{thebibliography}{99}
\bibitem [1]{Gomes} 
S. Gomes, P. Jacquinot, and M. Pisani, 2012. "The EAGLE. A model for policy analysis of macroeconomic interdependence in the euro area." Economic Modelling, Volume 29, Issue 5, Pages 1686-1714.
\bibitem [2]{Clancy} 
D. Clancy, P. Jacquinot, M. Lozej, 2016. "Government expenditure composition and fiscal policy spillovers in small open economies within a monetary union." Journal of Macroeconomics, Volume 48, Pages 305-326.
\bibitem [3]{EC} 
EC. 2019. "Report on Public Finances in EMU 2019", Part II - Performance of spending rules at EU and national level - a quantitative assessment, European Commission.
\bibitem [4]{Haroutunian} 
Haroutunian, S., K. Ba\'nkowski, S. Bischl, O. Bouabdallah, S. Hauptmeier, N. Leiner-Killinger, M. O'Connell, I. Arruga Oleaga, L. Abraham, and A. Trzcinska. 2024. "The path to the reformed EU fiscal framework: a monetary policy perspective." Occasional Paper Series, No 349, European Central Bank.
\end{thebibliography}


%----------------------------------------------------------------------------------------
%\end{multicols}
\end{document}
